\documentclass[10pt,a4paper]{amsart}
\usepackage[margin=19mm]{geometry}
\usepackage{amsmath,amssymb,mathtools}
\usepackage[scr=rsfs,cal=euler]{mathalfa}
\usepackage{xfrac}
\usepackage[unicode,hidelinks,bookmarks=false]{hyperref}
\newcommand{\ie}{i.e.\ }
\newcommand{\cf}{cf.\ }
\newcommand{\resp}{respectively\ }
\newcommand{\Subsection}{\subsection}
\providecommand{\nobreakdash}{\nobreak\hbox{-}}
\setlength{\emergencystretch}{2em}
\multlinegap0pt
\usepackage{bm}
\usepackage{bbm}
\usepackage{dsfont}
\usepackage{xspace}
\usepackage{cancel}
\usepackage{mathtools}
\usepackage{upgreek}
\usepackage{amsthm}
\makeatletter
\@ifundefined{thm}{\newtheorem{thm}{Thm}}{}
\@ifundefined{dcases}{\newtheorem{dcases}{Dcases}}{}
\@ifundefined{lemm}{\newtheorem{lemm}[thm]{Lemma}}{}
\makeatother
\newcommand*\abs[1]{\lvert#1\rvert}
\newcommand\Dcover[2]{\widetilde{#1^{#2}}}
\newcommand\etp[1]{\mathfrak{e}\left(#1\right)}
\newcommand\glpR{\mathrm{GL}_2^{+}(\mathbb{R})}
\newcommand\rmi{\mathrm{i}}
\newcommand\tbtmat[4]{\left(\begin{smallmatrix}{#1} & {#2} \\ {#3} & {#4}\end{smallmatrix}\right)}
\title[Dimension formulas for modular form spaces of rational weigh]{Dimension formulas for modular form spaces of rational weights, the classification of eta-quotient characters and an extension of Martin's theorem}
\author{Xiao-Jie Zhu}
\date{Source: arXiv:2408.00246v2 (CC BY 4.0)}
\begin{document}
\maketitle

\noindent\emph{Source and licence.} Dimension formulas for modular form spaces of rational weights, the classification of eta-quotient characters and an extension of Martin's theorem. Version arXiv:2408.00246v2 (CC BY 4.0), distributed under the Creative Commons Attribution 4.0 International licence, as recorded in the arXiv licence metadata for this exact version and shown on the abstract page https://arxiv.org/abs/2408.00246v2. Full rights notice: licenses/arxiv-2408.00246-CC-BY-4.0.txt.\par\smallskip

\noindent\textit{Editorial scope.} Counted unit: the Lemm whose source label is \texttt{lemm:neighborhood} in the article (source0.tex lines 3187--3201, source label \texttt{lemm:neighborhood}), delivered together with the article's own complete original proof of it (source0.tex lines 3202--3246, one \texttt{proof} environment of 45 lines). No mathematics is written, repaired or re-derived: statement and proof are the article's own text line for line. Required context retained uncounted: none. Every definition, display and label of those lines is retained unchanged. Omitted from this package and not redistributed: 1--3186, 3247--3429. Nothing inside a retained excerpt is omitted or altered: every retained body is delivered exactly as the article's lines are written, so no omission receipt of any kind accompanies this package. Citations: the retained excerpts carry no citation command at all, so no bibliography entry of the article is transcribed and no reference is redistributed. Reference adaptation: none was needed and none was made; every reference inside the delivered unit resolves inside it, so no unresolved reference or citation is produced. Printed slips kept literal: no printed word, symbol, hypothesis or number of the counted statement or of its proof was altered. Layout and class: the article's own class is \texttt{amsart} with packages including geometry, hyperref, amsmath, amsthm, amsfonts, amssymb, mathrsfs, mathtools, upgreek, longtable, booktabs; the wrapper is the lane's pinned \texttt{amsart} editorial wrapper at 10pt on A4, which restates the theorem-like environments the retained text needs and adds the article's own macro definitions that the retained text needs (\texttt{\textbackslash Dcover}, \texttt{\textbackslash abs}, \texttt{\textbackslash etp}, \texttt{\textbackslash glpR}, \texttt{\textbackslash rmi}, \texttt{\textbackslash tbtmat}), with extra wrapper packages (bm, bbm, dsfont, xspace, cancel, mathtools, upgreek). The delivered numbering is the unit's own. Assets: no retained span contains a figure environment, a graphics-inclusion command, a PDF-page inclusion, a table or a TikZ picture; no image file is copied into this package, the package declares no asset dependency and no image file is redistributed with it. Licence: version arXiv:2408.00246v2 of this article is distributed under the Creative Commons Attribution 4.0 International licence, as recorded in the arXiv licence metadata for this exact version and shown on the abstract page https://arxiv.org/abs/2408.00246v2. A full rights notice is delivered at licenses/arxiv-2408.00246-CC-BY-4.0.txt. Characters: every retained excerpt is pure ASCII, so the delivered engine, XeLaTeX with its default Latin Modern text font, reports no missing character.
\begin{lemm}
\label{lemm:neighborhood}
If $c\neq0$, or $c=0$ but $d>0$, then there exist $r_1,r_2>0$ such that for all $0<\delta_1<r_1$, we have
\begin{equation}
\label{eq:smallneighborhood1}
U(A;\delta_1,r_2)=\left\{(\gamma_1,\varepsilon_1\cdot\phi_{\gamma_1})\in\Dcover{\glpR}{D}\colon\lVert\gamma_1-\gamma\rVert<\delta_1,\,\varepsilon_1=\varepsilon\right\}.
\end{equation}
On the other hand, if $c=0$ and $d<0$, then there exist $r_1,r_2>0$ such that for all $0<\delta_1<r_1$, we have
\begin{multline}
\label{eq:smallneighborhood2}
U(A;\delta_1,r_2)=\left\{(\gamma_1,\varepsilon_1\cdot\phi_{\gamma_1})\in\Dcover{\glpR}{D}\colon\gamma_1=\tbtmat{a_1}{b_1}{c_1}{d_1},\,c_1\geq0,\,\lVert\gamma_1-\gamma\rVert<\delta_1,\,\varepsilon_1=\varepsilon\right\}\\
\cup\left\{(\gamma_1,\varepsilon_1\cdot\phi_{\gamma_1})\in\Dcover{\glpR}{D}\colon\gamma_1=\tbtmat{a_1}{b_1}{c_1}{d_1},\,c_1<0,\,\lVert\gamma_1-\gamma\rVert<\delta_1,\,\varepsilon_1=\varepsilon\cdot\etp{1/D}\right\},
\end{multline}
where the union is disjoint.
\end{lemm}

\begin{proof}
We only write down the proof for $c=0,d<0$, since the other case is simpler and similar. The first and second sets in the right-hand side of the desired identity are denoted by $V(A;\delta_1)_{+}$ and $V(A;\delta_1)_{-}$, respectively. We need to find $r_1$ and $r_2$ and to prove $U(A;\delta_1,r_2)=V(A;\delta_1)_{+}\cup V(A;\delta_1)_{-}$. Without loss of generality, we assume that $d'=d$, namely, $ad-bc=1$.

As a prerequisite, we define a function on a small disk centered at $d$:
\begin{equation*}
f(z)=\begin{dcases}
z^{\frac{1}{D}}=\exp\frac{1}{D}\log z &\text{ if }\Im z\geq0,\\
\etp{1/D}z^{\frac{1}{D}}=\exp\frac{1}{D}\left(\log z+2\uppi\rmi\right) &\text{ if }\Im z<0.
\end{dcases}
\end{equation*}
Then $f(z)$ is continuous on small disks not containing $0$. (Indeed, it is holomorphic.)

Set $l=\min_{0\leq i\neq j<D}\abs{\etp{i/D}-\etp{j/D}}$, and set $r_2=\frac{l}{2+l}\abs{d^{\frac{1}{D}}}$. By continuity of $f$ and boundedness of $\overline{B}$, we can find a $r_1>0$ with the property that for all $2\times2$ real matrices $\gamma_1=\tbtmat{a_1}{b_1}{c_1}{d_1}$ (not necessarily in $\glpR$)
\begin{equation}
\label{eq:r1withproperty}
\lVert\gamma_1-\gamma\rVert<r_1,\,c_1\geq0\Longrightarrow d_1\neq0,\,\abs{d^{\frac{1}{D}}-(c_1^\prime\tau+d_1^\prime)^{\frac{1}{D}}}<r_2\text{ for all }\tau\in\overline{B},
\end{equation}
and
\begin{equation}
\label{eq:r1withpropertyminus}
\lVert\gamma_1-\gamma\rVert<r_1,\,c_1<0\Longrightarrow d_1\neq0,\,\abs{d^{\frac{1}{D}}-\etp{1/D}(c_1^\prime\tau+d_1^\prime)^{\frac{1}{D}}}<r_2\text{ for all }\tau\in\overline{B}.
\end{equation}
We shall prove the above $r_1$ and $r_2$ suffice. Let $0<\delta_1<r_1$.

Let $(\gamma_1,\varepsilon_1\cdot\phi_{\gamma_1})\in V(A;\delta_1)_{+}$ be arbitrary. Then $c_1\geq0$, $\lVert\gamma_1-\gamma\rVert<\delta_1$, and $\varepsilon_1=\varepsilon$. This, together with \eqref{eq:r1withproperty}, implies $\lVert\varepsilon_1\cdot\phi_{\gamma_1}-\varepsilon\cdot\phi_{\gamma}\rVert<r_2$. Therefore, $V(A;\delta_1)_{+}\subseteq U(A;\delta_1,r_2)$. Now let $(\gamma_1,\varepsilon_1\cdot\phi_{\gamma_1})\in V(A;\delta_1)_{-}$ be arbitrary. Then $c_1<0$, $\lVert\gamma_1-\gamma\rVert<\delta_1$, and $\varepsilon_1=\varepsilon\cdot\etp{1/D}$. This, together with \eqref{eq:r1withpropertyminus}, implies $\lVert\varepsilon_1\cdot\phi_{\gamma_1}-\varepsilon\cdot\phi_{\gamma}\rVert<r_2$, so $V(A;\delta_1)_{-}\subseteq U(A;\delta_1,r_2)$.

Conversely, let $(\gamma_1,\varepsilon_1\cdot\phi_{\gamma_1})\in U(A;\delta_1,r_2)$ be arbitrary. Then $\lVert\gamma_1-\gamma\rVert<\delta_1$, and $\lVert\varepsilon_1\cdot\phi_{\gamma_1}-\varepsilon\cdot\phi_{\gamma}\rVert<r_2$. Set $\varepsilon_2:=\varepsilon$ if $c_1\geq0$ and $\varepsilon_2:=\etp{1/D}\varepsilon$ if $c_1<0$. The desired conclusion $(\gamma_1,\varepsilon_1\cdot\phi_{\gamma_1})\in V(A;\delta_1)_{+}\cup V(A;\delta_1)_{-}$ will follow if we can prove $\varepsilon_1=\varepsilon_2$. To this end, we note that
\begin{equation}
\label{eq:normInequality}
\abs{\varepsilon_1-\varepsilon_2}\cdot\lVert\phi_{\gamma_1}\rVert=\lVert\varepsilon_1\phi_{\gamma_1}-\varepsilon_2\phi_{\gamma_1}\rVert\leq\lVert\varepsilon_1\phi_{\gamma_1}-\varepsilon\phi_{\gamma}\rVert+\lVert\varepsilon\phi_{\gamma}-\varepsilon_2\phi_{\gamma_1}\rVert.
\end{equation}
We have $\lVert\varepsilon_1\phi_{\gamma_1}-\varepsilon\phi_{\gamma}\rVert<r_2$ by the assumption. By $\lVert\gamma_1-\gamma\rVert<\delta_1$, and by \eqref{eq:r1withproperty} or \eqref{eq:r1withpropertyminus} (according to $c_1\geq0$ or $c_1<0$, respectively), we have $\lVert\varepsilon\phi_{\gamma}-\varepsilon_2\phi_{\gamma_1}\rVert<r_2$. Inserting these into \eqref{eq:normInequality}, we have
\begin{equation}
\label{eq:normInequality2}
\abs{\varepsilon_1-\varepsilon_2}\cdot\lVert\phi_{\gamma_1}\rVert<2r_2.
\end{equation}
Again by \eqref{eq:r1withproperty} or \eqref{eq:r1withpropertyminus}, we have
$$\lVert\phi_{\gamma_1}\rVert=\lVert\varepsilon_2\phi_{\gamma_1}\rVert=\lVert\varepsilon_2\phi_{\gamma_1}-\varepsilon\phi_\gamma+\varepsilon\phi_\gamma\rVert\geq\lVert\phi_\gamma\rVert-\lVert\varepsilon_2\phi_{\gamma_1}-\varepsilon\phi_\gamma\rVert>\abs{d^{\frac{1}{D}}}-r_2.$$
Combining this and \eqref{eq:normInequality2} we find that
\begin{equation}
\label{eq:normInequality3}
\abs{\varepsilon_1-\varepsilon_2}<\frac{2r_2}{\abs{d^{\frac{1}{D}}}-r_2}=l,
\end{equation}
from which and the definition of $l$ we conclude that $\varepsilon_1=\varepsilon_2$. This finishes the whole proof.
\end{proof}

\end{document}
